\chapter{Rendering}
\label{ch:rendering}

\begin{figure}[H]
  \centering
  \resizebox{\textwidth}{!}{% The two render paths and where they meet: the views, F12, the display. Two rows so it reads at page width.
\begin{tikzpicture}[node distance=5mm and 6mm]
  \node[box, fill=blMint!12, text width=30mm] (scene) {The scene\\ \footnotesize evaluated meshes, materials, lights, world, camera};
  \node[box, right=of scene, yshift=12mm, text width=44mm] (raster) {Rasterizer (CPU)\\ \footnotesize \code{render/raster.*}: tile-binned, visibility buffer, deferred PBR, shadow maps, ids, outlines};
  \node[box, right=of raster, text width=40mm] (gpuv) {GPU views\\ \footnotesize \code{gpu::Raster} (Vulkan): the same formulas in GLSL};
  \node[box, right=of scene, yshift=-12mm, text width=44mm] (pt) {Path tracer (CPU)\\ \footnotesize \code{render/pathtracer.*}: BVH or Embree, NEE, MIS, GGX VNDF, denoisers};
  \node[box, right=of pt, text width=40mm] (gpupt) {GPU path tracing\\ \footnotesize \code{gpu\_kernel.h} (GLSL), OptiX; CPU and cards share one sample counter};
  \node[box, below=16mm of pt, xshift=-6mm, text width=40mm] (gi) {Lighting\\ \footnotesize lightmaps, probe volumes, voxel GI, reflection probes};
  \node[box, right=of gi, text width=44mm] (cm) {Colour management\\ \footnotesize Standard / Filmic / ACES, OpenColorIO views as LUTs, SIMD display encoding};
  \node[box, right=of cm, fill=blAmberFill!15, draw=blAmber, text width=44mm] (out) {The picture\\ \footnotesize Scene and Game views, Camera Preview, F12 and the Render window, camera filters, PNG / JPEG / HDR / EXR / MP4};
  \draw[flow] (scene) -- (raster);
  \draw[flow] (scene) -- (pt);
  \draw[flow, <->] (raster) -- (gpuv);
  \draw[flow, <->] (pt) -- (gpupt);
  \draw[flow] (gpuv.east) -- ++(4mm,0) |- (gi.east) ;
  \draw[flow] (gpupt.south) |- (gi.north east);
  \draw[flow] (gi) -- (cm);
  \draw[flow] (cm) -- (out);
  \node[note, below=4mm of gi.south west, anchor=north west, text width=150mm] {A view reuses its picture while a content hash of what it shows is unchanged (\adr{0009}); the Game view and the Camera Preview keep their last picture during a drag (\adr{0039}). Every kernel has a scalar reference the tests compare it with: the AVX2 display encoder, the 8-wide shading, the GPU views, the GPU kernel's samples (the same PCG seeds as the CPU's).};
\end{tikzpicture}
}
  \caption{The two render paths and where they meet. Every GPU path mirrors a CPU path with the same formulas, and the CPU path is the reference the tests compare it with.}
  \label{fig:render}
\end{figure}

\section{The rasterizer}
\code{render/raster.*} is a tile-binned, multithreaded software rasterizer with a z-buffer, an id buffer for picking and
outlines \cite{MarschnerShirley2021, Gregory2018vol2}. Its \code{Deferred} mode writes a visibility buffer and shades
once per pixel: GGX + Lambert \cite{Cook1982Reflectance, Schlick1994Inexpensive}, image-based lighting (SH9 irradiance
\cite{Ramamoorthi2001Efficient}, mip-filtered specular with the split-sum approximation \cite{Karis2013Real,
Hill2012}), sun shadow maps with PCF kept per caster set, area lights as emissive planes with an exact unshadowed
integral \cite{Heitz2018Combining}, toon and lilToon responses (\adr{0020}, \adr{0100}), lightmaps, probes and voxel GI
ranked by \code{light\_surface}. Eight same-triangle pixels of a plain material are lit at once with AVX2 (\adr{0035});
triangle setup rejects eight at a time; the scalar path is the reference and the tests compare them.

\section{The path tracer}
\code{render/pathtracer.*} is progressive: a two-level binned-SAH BVH with instancing \cite{Wald2007Fast,
Moller1997Fast, Duff2017Building} (or Embree's), next-event estimation of sun, point, spot, area and mesh lights with
multiple importance sampling, GGX sampling of visible normals \cite{Heitz2018Sampling, Walter2007Microfacet}, a thin
lens with bladed apertures \cite{Shirley1997Low}, stratified samples \cite{Cranley1976Randomization}, Russian
roulette, clamping, path guiding through Open PGL \cite{Ruppert2020Robust}, volumes (fog inside a closed mesh,
\adr{0092}), and two denoisers: OpenImageDenoise when it is there, else an edge-avoiding \`A-Trous filter
\cite{Dammertz2010Edge}. Transparent renders count coverage; filtered renders carry depth and id planes so camera
filters run on traced pictures too (\adr{0016}).

\section{The GPU}
\code{render/gpu\_kernel.h} is the megakernel in GLSL, a line-for-line port of the CPU's trace with the same PCG seeds,
so GPU and CPU samples are interchangeable; \code{gpu\_device.cpp} opens the Vulkan loader at run time, compiles the
kernel with shaderc (the SPIR-V cached on disk, pre-warmed by the build), uploads the scene as hashed sections and
renders sample batches (\adr{0013}, \adr{0014}). Devices open in the background and a render adopts its card when
ready (\adr{0021}); a worker per card keeps two submissions in flight (\adr{0052}); CPU threads and every GPU claim
sample indices from one atomic counter, the CPU stopping when the cards would finish sooner. The OptiX device
(\adr{0037}, \adr{0054}) runs the same kernel translated to CUDA with rays through \code{optixTrace}, its PTX committed
so the build needs no CUDA toolkit. The views are drawn by \code{gpu::Raster} (\adr{0036}, \adr{0061}): a shadow pass,
a G-buffer and full-screen lighting with the CPU's formulas, the Edit Mode cage, textures, toon maps, lightmaps and
probes read per visible pixel, overlays recorded in the world and replayed on the CPU when the card cannot; a window is
presented by the card one frame behind with only the changed tiles taken up. Each window's device is a preference and
the fall-back to the CPU is said once (a CPU badge in the Scene view's bar).

\section{Colour management}
HDR $\to$ view transform $\to$ 8-bit sRGB runs in \code{render/display.cpp} at 4 or 8 pixels a step with an exact sRGB
table (correctly rounded on 17 million floats in the test). The views are Standard, Filmic, Blendity's ACES curve,
Unity's ACES (Hill's fit of the RRT + ODT, \rnd{405}) \cite{Hill2016ACES, ACES2026}, and with OpenColorIO Blender's
config baked into $97^3$ tetrahedral LUTs over a log2 shaper, which the GPU views apply as the same table.
Light colour temperature follows the Planckian locus \cite{Kang2002Design}.

\section{Lighting}
Global illumination like Unity's (R-11, R-12): baked lightmaps (charts sized from world area, shelf-packed pages,
texels baked through the path tracer, saved beside the scene, \adr{0010}); probe volumes placed adaptively in bricks
of $4^3$ ambient cubes \cite{McTaggart2004} with validity, virtual offset and dilation (\adr{0012}); voxel GI
\cite{Thiedemann2011Voxel, Amanatides1987Fast, Dachsbacher2005, Hammersley1960} for one bounce of the sun and sky
occlusion in the rasterized views with nothing to bake (\adr{0011}), on the GPU as a compute kernel; realtime lightmaps
and live probes that gather through the voxel grid; reflection probes (\rnd{401}); mesh maps (normal, position,
curvature, occlusion, thickness) baked into the UV layout (\adr{0079}); texture bakes that turn a material into plain
images (\adr{0019}); and the world's Hosek--Wilkie sky \cite{Hosek2012Analytic}, HDRIs, sun disc, horizon band and fog.
A scene opened from Unity shows Unity's own lightmaps, probes, reflection probes, fog, skybox, post-processing and
BRDF weight as Unity shows them (\adr{0098}, \adr{0099}; rounds 400--405), measured against Unity's renders at fixed
stations (\code{tools/unity\_compare.py}).

\section{Camera filters}
Filters are components on a camera's object (\adr{0007}, \adr{0008}): an ordered stack of effects, each contributing
raster settings (vertex snap, affine UVs, texture rules, screen-door transparency: the PS1, N64, Saturn and DOS looks
\cite{psxspx}), an internal resolution, and image passes in linear or display light (colour grading in ASC CDL order
\cite{ASCCDL}, bloom \cite{Jimenez2014}, vignette, chromatic aberration, lens distortion \cite{Brown1966}, film grain,
CRT, pixelate, edge outline, sharpen, halftone, hatching, line art, paper, colour ramp, and the Anime and Manga
looks). The same stack runs on the Game view, the Camera Preview, F12 and camera sequences; the Scene view can look
through it. Filters look the same at every picture size: pixel-sized parameters scale with the height.
